\begin{table}[h!]
\centering
\footnotesize
\begin{tabular}{lccr}
\toprule
Attention dose & Mean STRAIN & Median STRAIN & Rusher-frames \\
\midrule
Singled (~1) & 0.1711 & 0.1576 & 632,492 \\
Doubled (~2) & 0.0127 & 0.0615 & 293,868 \\
3+ & -0.0672 & 0.0122 & 25,203 \\
\bottomrule
\end{tabular}
\caption{Descriptive dose--response of pressure on attention over the Phase-2.5 frame-level
assignments. Dose is the effective-blocker count $n_{j,t}=\sum_b\theta(b,j,t)$ on a rusher at a
frame. More attention coincides with \emph{less} strain --- a singled rusher averages
0.171, a doubled one 0.013 --- and the pooled slope of
STRAIN on dose is $-0.1017$ (cluster-robust SE 0.0019, clustered by
play). The suppressive effect of extra blockers dominates at the frame level, so unlike a season-level
``doubles chase winners'' correlation the descriptive sign is already correct. But the magnitude is
\emph{not} the causal effect. Part of the pooled slope is a between-rusher artifact (interior
rushers are both doubled more and generate less strain for positional reasons); restricting to within
a rusher (rusher fixed effect) the slope shrinks to $-0.0607$ (SE 0.0022).
What remains is biased \emph{toward zero} by selection: a defense commits the second blocker
precisely in the frames a rusher is threatening (a high-strain state), so the doubled frames are
over-sampled from moments whose strain would have been high anyway, masking part of the suppression.
The true causal effect of an added blocker is therefore \emph{more} negative than these descriptive
slopes; recovering its unbiased magnitude would require a causal design (for instance exploiting
stunts and twists as exogenous reallocations of attention), which we leave to future work.}
\label{tab:dose_response}
\end{table}
